\documentclass[10pt]{article}
\usepackage[margin=2.2cm]{geometry}
\usepackage{amsmath,amssymb,booktabs,microtype,xcolor,enumitem}
\usepackage[hidelinks]{hyperref}
\usepackage[small,compact]{titlesec}
\setlength{\parskip}{2pt}
\newcommand{\done}{\textcolor[rgb]{0,0.45,0}{\textbf{[DONE]}}}
\newcommand{\weak}{\textcolor[rgb]{0.8,0.5,0}{\textbf{[WEAK]}}}
\newcommand{\todo}{\textcolor[rgb]{0.75,0,0}{\textbf{[MISSING]}}}
\newcommand{\zslow}{z_{\mathrm{slow}}}
\newcommand{\zfast}{z_{\mathrm{fast}}}

\title{\vspace{-8mm}\large\bf HG-JEPA: State of the Evidence and a Plan Toward an
ICLR Submission\\ \normalsize (advisor briefing --- what is done, what is missing,
and whether it can clear the bar)}
\author{Jaechan Lee}
\date{July 25, 2026}

\begin{document}
\maketitle
\vspace{-6mm}

\section{The single message}
The project has drifted through several framings (separation guarantee
$\to$ assignment mechanism $\to$ negative anomaly pilot). For the paper we
commit to \textbf{one} thesis and make every result either support it or
scope it:

\begin{quote}
\textbf{The prediction horizon is a free supervision axis.} In any
latent-target predictive learner, gating a designated latent block out of the
predictor at long horizons provably \emph{assigns} the slow generative
factors to the remaining block --- with no labels --- while a narrow
bottleneck and a decorrelation penalty handle \emph{exclusion} of fast
factors, which we prove no objective of this family can enforce by itself. We
characterize exactly when this works (a genuine timescale gap) and when it
does not (XJTU bearings, non-novelty drift), and validate the mechanism
across two objectives (JEPA regression, CPC-InfoNCE) and three real datasets.
\end{quote}

This is a \emph{method paper whose claims are scoped by a mechanism}, the
standard shape of a defensible ICLR submission. It is not a retreat from the
method framing: the deliverable is still ``a label-free separator you can add
to any latent predictive model,'' but every claim is stated at the strength
the evidence supports --- assignment is structural (Proposition 1), exclusion
is linear-subspace-level (Proposition 2 + MLP/MINE measurements), and
applicability is diagnosable before training ($T_{\mathrm{ac}}$ screening).
What we explicitly stop claiming: information-theoretic removal of fast
factors, anomaly attribution as a working payoff, and any performance
guarantee without a timescale gap.

\paragraph{How the claim evolved (one paragraph, for the paper's own honesty
and for this meeting).} We began with ``the gate separates slow from fast''
(H1, strong form). Experiments forced three revisions: (i) the gate alone
separates nothing --- assignment vs.\ exclusion had to be split (now
Propositions~1/2); (ii) linear probes overstated exclusion --- an MLP
recovers the fast factor from $\zslow$ at $R^2\,0.45$ (vs.\ $0.19$ linear),
MINE agrees ($0.95$ vs.\ $1.21$ nats ungated), so the separation claim is
linear-subspace; (iii) the flagship payoff (anomaly attribution) failed in
its contextual half (detection AUROC $\approx 0.50$) because the encoder
projects unseen dynamics onto normal slow codes. Each revision shrank the
claim and strengthened its defensibility; the paper should present this as
mechanism characterization, not as a wandering narrative.

\section{Critical assessment: can this clear ICLR?}
Frank verdict first: \textbf{as it stands today this is a solid TMLR paper
and a borderline-reject at ICLR.} TMLR's bar is ``claims supported by
evidence,'' which we now meet after the honesty passes. ICLR additionally
demands either breadth of empirical impact or depth of theory; we currently
have two positive real datasets, no external SOTA baselines, and two
propositions that are useful but not an identifiability theorem.

\textbf{What we have that reviewers will like.} A clean, falsifiable
mechanism with matching theory (assignment provable; exclusion provably
unidentifiable --- and the measured nonlinear leak lands exactly on that
boundary); a leak-free evaluation protocol (group splits, absolute scores,
RankMe) that is stricter than most of the disentanglement literature; honest
negative results (XJTU, contextual anomalies) that double as evidence the
$T_{\mathrm{ac}}$ screening predicts applicability; cross-objective
generality (the CPC-InfoNCE cell separates \emph{best}: regime $0.81$, leak
$0.05/0.00$); and an unusual-for-the-field control showing post-hoc unmixing
(SFA/ICA/PCA) cannot substitute for training-time gating.

\textbf{The five gaps that decide acceptance} (each maps to a work item in
\S\ref{sec:plan}):
\begin{enumerate}[leftmargin=1.4em,itemsep=1pt,topsep=2pt]
\item \emph{Dataset breadth.} Two positive real datasets (HAPT, PTB-XL) is
below the bar. Add Sleep-EDF (EEG sleep stages; large timescale gap, subject
split) and C-MAPSS (turbofan degradation; unit split) with the
$T_{\mathrm{ac}}$ gap screening \emph{pre-registered}: predict
separable/not-separable before training, then verify. This turns the XJTU
negative into predictive power.
\item \emph{The downstream cost must be owned and answered.} On HAPT,
gate$+$dcor training costs the full embedding $0.80\!\to\!0.76$ activity F1,
and deploying $\zslow$ alone costs $0.80\!\to\!0.71$ vs.\ an ungated model
--- currently under-reported. The paper must present the compensation
explicitly: selective robustness (the HAPT result already holds on unseen
subjects; extend to corruption/domain-shift stress tests where the full
embedding degrades and $\zslow$ does not), and the synthetic case where
$\zslow$ \emph{beats} the full embedding ($0.80$ vs.\ $0.74$) because fast
dimensions are noise for the slow task.
\item \emph{Label-free model selection} (the Locatello objection, our
reviewer \#3). $\tau$ is already set label-free from $T_{\mathrm{ac}}$;
$d_{\mathrm{slow}}$ and $\lambda$ are not. A retrospective validation on a
9-config $\tau\times d_{\mathrm{slow}}\times\lambda$ grid (2 seeds) is now
in: ranking by unsupervised health metrics (per-block RankMe, cross-block
correlation, temporal-roughness gap) correlates significantly with the
labeled separation score (Spearman $0.68$, $p{=}0.04$) and correctly
\emph{rejects} every broken configuration (no-gate, gate-only,
$\lambda{=}16$ degeneracy) --- but among healthy configurations it picks
$\lambda{=}1$ where labels prefer $d_{\mathrm{slow}}{=}8,\lambda{=}4$
(top-1 regret $0.27$ on a $0.66$ scale). Honest position for the paper:
the criterion is a validated \emph{catastrophe filter}, not a tuner; we
therefore fix one global default ($d_{\mathrm{slow}}{=}16,\lambda{=}4$)
across all datasets, use the filter as the only per-dataset decision, and
state that fine selection still requires labels (the residual Locatello
concession). Real-data replication of the filter \todo.
\item \emph{External baselines.} SlowVAE is refuted
(loses the regime entirely, $0.63$/full $0.54$), but we have no
self-supervised SOTA representation baseline. Add TS2Vec and a masked
reconstruction baseline under matched budgets on all real datasets ---
reviewers will not accept internal ablations alone.
\item \emph{Close or scope the nonlinear leak.} One attempt at a nonlinear
dependence penalty (HSIC between blocks, drop-in for dcor) to see whether the
$0.45$ MLP leak closes. If it does not, the paper claims linear-subspace
separation explicitly, with Proposition~2 as the reason this is a boundary
rather than a bug.
\end{enumerate}

\textbf{Three risks that can kill the submission.} (a) Sleep-EDF/C-MAPSS
fail to separate despite passing the gap screen $\to$ the method looks
synthetic$+$HAPT-specific; fallback is the TMLR analysis framing. (b) The
unsupervised selection criterion fails $\to$ the Locatello objection stands
unanswered; we would then fix $d_{\mathrm{slow}}{=}16$, $\lambda{=}4$
\emph{globally across all datasets} and show results are acceptable without
any tuning, a weaker but honest defense. (c) Reviewers reject
``linear-subspace separation'' as too weak a deliverable; mitigation is risk
(a)'s breadth plus the robustness payoff --- the claim must earn its keep in
a use case, not in metric aesthetics. \textbf{Go/no-go:} run Sleep-EDF first
(fastest, highest information). If it reproduces the HAPT pattern, commit to
ICLR; if not, submit the mechanism/conditions paper to TMLR.

\section{The paper, section by section}
\label{sec:paper}
For each section: what goes in it, what exists, what is missing.

\subsection{Introduction}
\textbf{Content.} (1) Practitioners need the slow state (health, diagnosis,
activity) separated from fast dynamics, not merely encoded --- motivating
uses: selective domain robustness and interpretable health trajectories
(anomaly attribution is demoted to a discussed-and-tested limitation).
(2) The one-sentence idea: the prediction horizon, a free knob in every
predictive SSL objective, is itself a supervision signal for
timescale factoring. (3) Claims list with exact strength (assignment
provable / exclusion linear / applicability diagnosable).
\textbf{Status:} motivation prose exists in the 2-page report \done;
horizon-as-supervision framing is currently buried in the method section
\weak; the claims list needs rewriting at the new strength \todo.

\subsection{Related work}
\textbf{Content and the references that were missing from the 2-page
report.} Four lines, one paragraph each, each ending with the delta:
(1) \emph{Slowness and predictive information}: SFA
\cite{wiskott2002slow}; predictive information \cite{bialek2001predictability};
past--future information bottleneck \cite{creutzig2009past}. Delta: slowness
of a feature $\neq$ long-horizon predictive value; our Tab.~1 control makes
this empirical.
(2) \emph{Identifiability}: TCL \cite{hyvarinen2016unsupervised}, PCL
\cite{hyvarinen2017nonlinear}, SlowVAE \cite{klindt2021towards}, the
impossibility result \cite{locatello2019challenging}. Delta: they get
theorems under generative assumptions; we get a weaker but
architecture-level guarantee (assignment) that composes with any
latent-predictive objective.
(3) \emph{Sequential VAEs / content--dynamics}: DSAE \cite{hsu2017unsupervised},
C-DSVAE \cite{bai2021contrastively}. Delta: static content vs.\ our
slowly-varying factor; reconstruction is local (our SlowVAE result: the
regime never enters the code, $0.54$ full-embedding accuracy).
(4) \emph{Latent-predictive SSL}: CPC \cite{oord2018representation}, I-JEPA
\cite{assran2023self}, data2vec \cite{baevski2022data2vec}, HEPA
\cite{hepa2026}, VICReg \cite{bardes2022vicreg} (for the decorrelation
penalty's lineage), RankMe \cite{garrido2023rankme}, MINE
\cite{belghazi2018mine}, DCI \cite{eastwood2018framework}, MIG
\cite{chen2018isolating}, TS2Vec \cite{yue2022ts2vec}.
Delta: these learn multi-scale content but do not factor it; we add the
factoring axis.
\textbf{Status:} positioning paragraph exists \done; full related-work
section with per-item deltas and complete citations \todo\ (this document's
bibliography is the seed).

\subsection{Method}
\textbf{Content --- must be fully self-contained (this was a gap in the
2-page report).}
\emph{Backbone.} Causal transformer encoder (4 layers, $d_{\mathrm{model}}
{=}96$, $\sim$0.5M params) over patches of $P{=}8$ steps, window $L{=}256$
patches; embedding $z_t\in\mathbb R^{64}$ split as
$[\zslow(16);\zfast(48)]$; EMA target encoder (0.996); predictor
$P(z,\Delta)$ conditioned on horizon $\Delta\in\{1,4,16,64,128\}$ patches;
objectives: latent $L_2$ regression (JEPA family) and, in the new cell,
InfoNCE against EMA embeddings with in-batch negatives (CPC family).
\emph{The gate.} $g(\Delta)=\sigma((\tau-\Delta)/w)$ multiplies $\zfast$ at
the predictor input; asymmetric by design (short horizons see both blocks,
because fast dynamics are conditioned on slow state).
\emph{Decorrelation.} Squared cross-covariance between blocks, weight
$\lambda$; VICReg-lineage.
\emph{Label-free $\tau$.} $\tau=c\,T_{\mathrm{ac}}$, $c\!\approx\!2.5$, from
the signal's autocorrelation decay --- also the pre-training applicability
screen.
\textbf{Status:} all implemented and stable across 3 seeds \done; the
description above must be written out with a figure (architecture + gate
schedule) \todo; HSIC variant \todo.

\subsection{Theory}
\textbf{Content.} Proposition 1 (inclusion): under a timescale gap
($Y_\Delta\perp X_{\le t}\mid s_t$ for $\Delta\ge\Delta_0$) and a closed
gate, any population minimizer makes $\zslow$ mean-sufficient for the far
future; with affinely independent regime-conditional target means, the
slow-state posterior is a function of $\zslow$. Proposition 2 (no
exclusion): a Borel-injection construction shows fast-information membership
in $\zslow$ is not identified by the objective at any gate schedule; dcor
removes linearly readable duplication only. Corollary: the measured
nonlinear leak is the identified/unidentified boundary, not an
implementation failure.
\textbf{Status:} both propositions written with proofs in the current report
\done; a full identifiability theorem (uniqueness up to symmetry, TCL-style)
is explicitly out of scope for September --- state as open \done\ (as a
statement); anonymous-ready notation cleanup \todo.

\subsection{Experiments}
\textbf{Datasets (each with the one-line reason it is in the paper).}
\emph{Synthetic two-timescale oscillator}: 3-state Markov regime (dwell 3000
steps) sets carrier frequency $\pm5\%$; fast OU factor (lifetime 50 steps)
jitters frequency comparably; amplitudes equalized so no short-window
statistic reveals the regime --- ground truth for mechanism claims.
\emph{HAPT} \cite{reyes2016transition}: smartphone IMU, slow $=$ activity
($\sim$12\,s dwell), fast $=$ gait cycle; subject-level split --- the
domain-robustness testbed. \emph{PTB-XL} \cite{wagner2020ptb}: 12-lead ECG,
slow $=$ diagnosis (static-content limitation flagged), fast $=$ beat
morphology; patient split. \emph{XJTU-SY} \cite{wang2020xjtu}: bearing
vibration, the \emph{pre-registered negative} --- carrier and envelope share
one timescale, gate correctly does nothing. Planned: \emph{Sleep-EDF}
\cite{kemp2000analysis} and \emph{C-MAPSS} \cite{saxena2008damage}.
\textbf{Protocol (non-negotiable, in the paper verbatim).} Probes are fit on
training groups and evaluated on held-out groups (subject/patient/bearing;
contiguous time split on synthetic) over disjoint windows; absolute scores
only (a ratio once hid a collapse that absolute scores + RankMe exposed);
$n{=}3$ seeds; all variants share one backbone and budget.
\textbf{Result inventory and status:}
\begin{itemize}[leftmargin=1.3em,itemsep=0pt,topsep=1pt]
\item F1 assignment + unmixing control (SFA/ICA/PCA cannot substitute) \done
\item F2 objective generality: raw-AR fails; JEPA $0.19/0.27$; CPC-InfoNCE
$0.05/0.00$ leak at regime $0.81$ \done
\item F3 gate-alone/dcor-alone factorization (matches Prop.~2) \done
\item F4 width sweep: absolute block size governs exclusion; scales to
$d_z{=}256$ \done
\item Nonlinear honesty: MLP $0.45$, MINE $0.95$ nats \done\ (HSIC response
\todo)
\item F5 real data: HAPT/PTB-XL positive, XJTU negative-as-predicted \done;
Sleep-EDF, C-MAPSS \todo
\item F6 competitors: SlowVAE refuted \done; TS2Vec, masked-recon baselines
\todo
\item DCI/MIG reported with the axis-metric caveat \done
\item Downstream cost/payoff table ($\zslow$ vs.\ ungated-full, all datasets,
plus robustness stress test) \weak\ (numbers exist for cost; robustness
stress test beyond unseen-subject \todo)
\item Label-free selection: retrospective grid validation \weak\ (running;
real-data replication \todo)
\item Anomaly attribution: point works ($0.93$--$0.97$), contextual at
chance --- reported as a scoped negative with diagnosis \done
\end{itemize}

\subsection{Conclusion / limitations}
\textbf{Content.} Restate the claim at its proven strength; limitations
carried verbatim from the current report (soft exclusion, two-timescale
assumption, $\tau$'s ACF assumption, patch-rate bound, dcor data-dependence,
validation-needs-labels); the anomaly negative and what would fix it
(novelty-sensitive machinery). \textbf{Status:} all content exists \done;
consolidation \todo.

\section{Workplan to the September deadline}
\label{sec:plan}
Ordered by information-per-day; each item names its kill criterion.
\begin{itemize}[leftmargin=1.3em,itemsep=1pt,topsep=2pt]
\item \textbf{W1 (Jul 28--Aug 1): Sleep-EDF end-to-end} (prep, screen,
train, probe). Go/no-go for the ICLR framing.
\item \textbf{W2: C-MAPSS} + downstream robustness stress tests
(corruption/shift) on HAPT$+$Sleep-EDF.
\item \textbf{W3: label-free selection} finish synthetic grid, replicate on
HAPT; HSIC penalty experiment.
\item \textbf{W4: external baselines} (TS2Vec, masked recon) on all real
datasets, matched budgets.
\item \textbf{Sep W1--2: writing} to this document's outline; internal
red-team pass against the 10-point reviewer critique we already answered
once.
\item \textbf{Slack:} the anomaly follow-up (density model on $\zslow$) is
explicitly \emph{out} of the September scope.
\end{itemize}

\begin{thebibliography}{99}\small\itemsep0pt
\bibitem{wiskott2002slow} L. Wiskott, T. Sejnowski. Slow feature analysis:
unsupervised learning of invariances. \emph{Neural Computation}, 2002.
\bibitem{bialek2001predictability} W. Bialek, I. Nemenman, N. Tishby.
Predictability, complexity, and learning. \emph{Neural Computation}, 2001.
\bibitem{creutzig2009past} F. Creutzig, A. Globerson, N. Tishby.
Past-future information bottleneck in dynamical systems. \emph{Phys.\ Rev.\
E}, 2009.
\bibitem{hyvarinen2016unsupervised} A. Hyv\"arinen, H. Morioka. Unsupervised
feature extraction by time-contrastive learning and nonlinear ICA.
\emph{NeurIPS}, 2016.
\bibitem{hyvarinen2017nonlinear} A. Hyv\"arinen, H. Morioka. Nonlinear ICA
of temporally dependent stationary sources. \emph{AISTATS}, 2017.
\bibitem{klindt2021towards} D. Klindt et al. Towards nonlinear
disentanglement in natural data with temporal sparse coding. \emph{ICLR},
2021.
\bibitem{locatello2019challenging} F. Locatello et al. Challenging common
assumptions in the unsupervised learning of disentangled representations.
\emph{ICML}, 2019.
\bibitem{hsu2017unsupervised} W.-N. Hsu, Y. Zhang, J. Glass. Unsupervised
learning of disentangled and interpretable representations from sequential
data. \emph{NeurIPS}, 2017.
\bibitem{bai2021contrastively} J. Bai, W. Wang, C. Gomes. Contrastively
disentangled sequential variational autoencoder. \emph{NeurIPS}, 2021.
\bibitem{oord2018representation} A. van den Oord, Y. Li, O. Vinyals.
Representation learning with contrastive predictive coding. arXiv:1807.03748,
2018.
\bibitem{assran2023self} M. Assran et al. Self-supervised learning from
images with a joint-embedding predictive architecture. \emph{CVPR}, 2023.
\bibitem{baevski2022data2vec} A. Baevski et al. data2vec: a general framework
for self-supervised learning in speech, vision and language. \emph{ICML},
2022.
\bibitem{hepa2026} J. Petersen, G.-A. Lombardi, R. Maggioni, C. Mazzoleni, F.
Martelli, and P. Petersen. HEPA: a self-supervised horizon-conditioned event
predictive architecture for time series. \emph{arXiv:2605.11130}, 2026.
\bibitem{bardes2022vicreg} A. Bardes, J. Ponce, Y. LeCun. VICReg:
variance-invariance-covariance regularization for self-supervised learning.
\emph{ICLR}, 2022.
\bibitem{garrido2023rankme} Q. Garrido et al. RankMe: assessing the
downstream performance of pretrained self-supervised representations by
their rank. \emph{ICML}, 2023.
\bibitem{belghazi2018mine} M. Belghazi et al. Mutual information neural
estimation. \emph{ICML}, 2018.
\bibitem{eastwood2018framework} C. Eastwood, C. Williams. A framework for
the quantitative evaluation of disentangled representations. \emph{ICLR},
2018.
\bibitem{chen2018isolating} R.T.Q. Chen et al. Isolating sources of
disentanglement in variational autoencoders. \emph{NeurIPS}, 2018.
\bibitem{yue2022ts2vec} Z. Yue et al. TS2Vec: towards universal
representation of time series. \emph{AAAI}, 2022.
\bibitem{reyes2016transition} J.-L. Reyes-Ortiz et al. Transition-aware
human activity recognition using smartphones. \emph{Neurocomputing}, 2016.
\bibitem{wagner2020ptb} P. Wagner et al. PTB-XL, a large publicly available
electrocardiography dataset. \emph{Scientific Data}, 2020.
\bibitem{wang2020xjtu} B. Wang et al. A hybrid prognostics approach for
estimating remaining useful life of rolling element bearings (XJTU-SY
datasets). \emph{IEEE Trans.\ Reliability}, 2020.
\bibitem{kemp2000analysis} B. Kemp et al. Analysis of a sleep-dependent
neuronal feedback loop: the slow-wave microcontinuity of the EEG. \emph{IEEE
Trans.\ Biomed.\ Eng.}, 2000. (Sleep-EDF)
\bibitem{saxena2008damage} A. Saxena et al. Damage propagation modeling for
aircraft engine run-to-failure simulation. \emph{PHM}, 2008. (C-MAPSS)
\end{thebibliography}

\end{document}
